\section{A semantic basis for framework comparison}\label{sec:C22:axes}
Define a framework profile by five sets: states it represents, transitions it controls, effects it mediates, observations it exports, and obligations it leaves to the application. This representation makes omissions visible. An SDK can validate a tool argument but leave external idempotency unspecified. A runtime can persist graph state but leave document authorization to a connector. A managed service can maintain sessions while exposing only partial execution traces. None of these limitations is inherently disqualifying; they determine the boundary of a defensible claim.

A useful comparison matrix distinguishes documented support from verified behavior. ``Persistence'' should be decomposed into conversation continuity, process-restart recovery, durable external artifacts, semantic memory, and cross-version migration. ``Human review'' should identify where execution pauses, what the reviewer sees, and whether authority is rechecked after resumption. ``Observability'' should identify whether traces include full tool arguments, truncation indicators, source identity, and asynchronous work. A single Boolean column loses these distinctions.

The unit of versioning is similarly broader than a package number. An application may combine a pinned SDK, moving model alias, remote tool schemas, hosted runtime, and changing connector permissions. Recording the SDK release alone does not reproduce the system. Define a compatibility vector whose components include each mutable boundary and attach it to evaluation traces. Unknown service revisions should remain unknown rather than be inferred from a documentation date.

More precisely, a configured deployment has a manifest $\xi=(\nu,\theta,\delta,\tau,\beta,\alpha)$: application/framework revision $\nu$, model identity $\theta$, deployment decoder $\delta$, tool-schema/adapter revision $\tau$, persistence backend/schema $\beta$, and current authority epoch $\alpha$. These components need not share a release cycle. A saved run additionally requires its realized context and external-effect records; $\xi$ identifies configuration rather than reconstructing those data. Framework comparison holds the relevant components fixed or declares their joint change as the treatment. Otherwise a newer SDK coupled to a stronger model is not an estimate of an SDK effect.

For feature equivalence, require observable contracts. Two frameworks both exposing a method called resume are not equivalent if one restores a model-visible history and another resumes an event-sourced workflow. The first may repeat tool effects; the second may skip completed activities but depend on deterministic application code. A conformance workload should therefore include an interruption after external acceptance and before local commitment, rather than test only that a suspended conversation can receive another message.

The practical implication is to select a framework against the failure modes of the intended workload. A short extraction service may benefit most from typed output and dependency injection. A multi-hour engineering task needs workspace state, cancellation, and reconciliation. A regulated business action needs principal identity and authoritative permission checks. Breadth of advertised features cannot replace the task-specific evidence needed to compare these roles.

